% !PRINTED 590
\weberpage{590}
%<<CONT
zerlegen, die zueinander relativ prim sind, und man kann dann
zwei ungerade ganze Zahlen $\alpha_1, \alpha_2$ in $\Omega$ bestimmen, die zueinander
relativ prim sind, von denen die eine durch $\mathfrak{a}_1$, die andere
durch $\mathfrak{a}_2$ teilbar ist. Es ist dann $\eta\alpha_1\alpha_2$ eine ganze Zahl in $\Omega$
und man kann nach Bd.\ II, \S\,166, (4) zwei ganze Zahlen $\varrho_1, \varrho_2$
so bestimmen, da\ss{}
\[
\frac{\nu\,\alpha_1\alpha_2}{\mu} = \varrho_1\alpha_2 + \varrho_2\alpha_1
\]
und daher
\[
\boldsymbol{s}\!\left(\frac{2\nu}{\mu}\right) = \boldsymbol{s}\!\left(\frac{2\varrho_1}{\alpha_1} + \frac{2\varrho_2}{\alpha_2}\right)
\]
wird. Dies ist aber von der Form (25). Wir bekommen daraus
den Satz:

4. \gesperrt{Enth\"alt eine gebrochene Zahl $\eta$ in $\Omega$ mit ungeradem
Nenner zwei oder mehr verschiedene Primfaktoren im Nenner, so
ist $\boldsymbol{s}(\eta)$ eine Einheit.}

\begin{center}
\textbf{\S\ 158. \gesperrt{Zerlegung der Funktion $\boldsymbol{A}(x)$.}}
\end{center}

Es seien jetzt
\weq{1}{\begin{aligned} \mu &= y + x\sqrt{\Delta}, \\ \mu_1 &= y_1 + x_1\sqrt{\Delta} \end{aligned}}
zwei ungerade ganze Zahlen in $\Omega$ von der Form \S\,157, (6). Es
sei $\mathfrak{m}$ der gr\"o\ss te gemeinschaftliche Idealteiler von $\mu$ und $\mu_1$.
Setzt man also
\weq{2}{\begin{aligned} \mu &= \mathfrak{m}\mathfrak{a}, \\ \mu_1 &= \mathfrak{m}\mathfrak{a}_1 \end{aligned}}
so sind $\mathfrak{a}$ und $\mathfrak{a}_1$ zwei \gesperrt{\"aquivalente Ideale} ohne gemeinsamen
Teiler.

Die den beiden Zahlen $\mu, \mu_1$ entsprechenden Funktionen $\boldsymbol{A}(x)$
bezeichnen wir mit
\[
\boldsymbol{A}_\mu(x), \ \boldsymbol{A}_{\mu_1}(x),
\]
und fragen, welche gemeinschaftliche Wurzeln diese Funktionen
haben, wann also die Gleichung:
\weq{3}{\boldsymbol{s}\!\left(\frac{2\varrho}{\mu}\right) = \boldsymbol{s}\!\left(\frac{2\varrho_1}{\mu_1}\right)}
erf\"ullt sein kann, wenn $\varrho$ nach dem Modul $\mu$, $\varrho_1$ nach dem Modul $\mu_1$
genommen ist.
